% META: {"id":"1SPE-DERLOCAL-CO-037","chapitre":"1SPE-DERIVATION-LOCAL","type_objet":"corrige","exercice_id":"1SPE-DERLOCAL-EX-037","status":"approved","origin":"generated","status_history":["generated","audited","accepted","approved"],"release_acceptance":"RELEASE_OWNER_BATCH_ACCEPTANCE","acceptance_closure_digest":"sha256:9b3ccf9a81c5520fbb7e03b7b2d3e7bf2057a02834b6d908bf3be5f49f3b553f","acceptance_receipt_digest":"sha256:4fad86f0282e0fa4e0419cca1ad6379ae7af5a280c04a4a90880f90a1c044242"}
% BEGIN-VERIFY
% from sympy import symbols
% x = symbols('x')
% f = x**2
% assert f.subs(x, 3) == 9
% fp = 2*x
% assert fp.subs(x, 3) == 6
% tangente = 6*(x - 3) + 9
% assert tangente.expand() == 6*x - 9
% END-VERIFY
\begin{corrige}{1SPE-DERLOCAL-EX-037}
On rappelle que l'équation de la tangente à la courbe de $f$ au point d'abscisse $a$ est :
\[
y = f'(a)(x - a) + f(a)
\]

On calcule $f(3)$ :
\[
f(3) = 3^2 = 9
\]

On admet que $f'(x) = 2x$ (dérivée de $x^2$), donc :
\[
f'(3) = 2 \times 3 = 6
\]

L'équation de la tangente au point d'abscisse $3$ est :
\[
y = 6(x - 3) + 9 = 6x - 18 + 9
\]

\[
\boxed{y = 6x - 9}
\]
\end{corrige}
